\section{Related work}
\label{sec:related}

\paragraph{Prompt injection in lending.} PIBench~\cite{roelofs2026pibench} is a hand-authored benchmark of 78 cases (53 injections) against an underwriting agent, evaluated on 16 hosted frontier models with and without untrusted-content tagging. It evaluates the LLM as decider and reports that tagging raises average defense success from 93.2\% to 97.7\%. It is a vendor web page and has not been peer reviewed. The architecture comparison in this paper, the scorecard-grounded definition of harm and the open-weight setting are outside its scope.

\paragraph{Document-borne injection in hiring.} A measurement over about 200,000 real resumes finds hidden injections and reports that over 90\% of injected prompts carry data and no instruction~\cite{zhang2026resume}. ResumeShield evaluates filtering plus an application-layer data fence on 104 synthetic resumes and lists falsified visible claims as outside its scope~\cite{barach2026}. A study from a recruitment marketplace tests ten instruction-style injection strings against an LLM that labels resume-to-job matches, with 24 prompting variants on hosted OpenAI models~\cite{akdemir2025}. None of these places a deterministic decider behind the extractor or works in lending.

\paragraph{Architectural defenses.} The design-pattern analysis describes action-selector, plan-then-execute, map-reduce, dual-LLM, code-then-execute and context-minimization patterns and applies them to ten case studies, including resume screening, without measuring attack success~\cite{beurerkellner2025}. CaMeL enforces control-flow and data-flow policies for tool-using agents and lists as an explicit non-goal attacks that change content without changing data flow~\cite{debenedetti2025camel}. Provenance and corroboration checks on individual fields appear in agent settings~\cite{fan2026pact,pandey2026}. Architecture B corresponds to the quarantined-extractor half of the dual-LLM pattern with a trained scorecard in place of a privileged LLM, and C adds checks of the kind these works describe.

\paragraph{Prompt-level defenses.} Delimiting or tagging untrusted content~\cite{hines2024} and training models to rank instructions by privilege~\cite{wallace2024} are the sources of our prompt-defended comparison point. Adaptive attacks bypass 12 recent defenses of this family~\cite{nasr2025}. We test a fixed attack set, so our estimates for prompt-level defenses are optimistic by construction.

\paragraph{Feature manipulation.} Agents who change their features to obtain a favorable classification are the subject of strategic classification~\cite{hardt2016}. The falsification variants here are a lending instance of that problem passing through an LLM.
